\documentclass[../main.tex]{subfiles}
\begin{document}
%==========================================TIÊU ĐỀ TẬP========================================
\begin{table}[h]
    \begin{adjustwidth}{-1cm}{}
        \begin{tabular}{>{\centering\arraybackslash}p{6cm}|>{\raggedright\arraybackslash}p{14cm}}
            \multirow{5}{6cm}
            % Cột bên trái: ÉPISODE và số tập
            {\\[-40pt]
                \flaregothic\fontsize{20pt}{20pt}\selectfont \centering
                \color{eptype}HISTOIRE\color{black}\\
                \brian\fontsize{72pt}{72pt}\selectfont
                \color{epnum}1
                \\[18pt]
                % \flaregothic\fontsize{14pt}{14pt}\selectfont
                % \color{epattr}BIENVENUE, \uline{\color{Banpire} Mel} !
                % \flaregothic\fontsize{18pt}{18pt}\selectfont
                % \color{epattr}ÉPISODE PERDU
                \color{black}
            }
             &
            % Dòng thứ nhất: tên tập tiếng Nhật
            {\fotskip\fontsize{18pt}{18pt}\selectfont \ruby{誕}{たん}\ruby{生}{じょう}\ruby{日}{び}\ruby{作}{さく}\ruby{戦}{せん}、\ruby{始}{し}\ruby{動}{どう}！} \\[6pt]
            % Dòng thứ hai: tên tập tiếng Anh
             & {\cafeteria\fontsize{18pt}{18pt}\selectfont Birthday Operation: Begin!}                                                       \\[4pt]
            % Dòng thứ ba: tên tập tiếng Việt
             & {\vnmsans\fontsize{18pt}{18pt}\selectfont Mật vụ sinh nhật bắt đầu!}                                                          \\[8pt]
            % Dòng thứ tư: Nhân vật xuất hiện
             & {\sen Track Used: Birthday Party (\ruby{誕}{たん}\ruby{生}{じょう}パーティー)}
        \end{tabular}
    \end{adjustwidth}
\end{table}
%===========================================NỘI DUNG==========================================
\color{black}

Ngày Mười Bảy tháng Bảy.

Một ngày bình thường như bao ngày khác. Nhưng chỉ có một điều khác biệt khiến mọi thứ trở nên kỳ lạ: một sự vắng mặt bất thường đã làm xáo trộn không khí thường ngày.

Buổi sáng hôm đó, văn phòng \hll\ vẫn như mọi khi: các cô gái đến đông đủ, trò chuyện rôm rả về những chủ đề không đầu không cuối. Một nhóm bàn về buổi phát trực tiếp tối qua, nhóm khác tám chuyện về quán ăn vừa mở gần đây, trong khi vài người đang gặm dở bữa sáng. Âm thanh cười nói, tiếng gõ bàn phím, và cả những tiếng ``phụt!'' khi ai đó cười sặc lên vì một chuyện hài hước nào đó vang vọng khắp căn phòng.

Mọi thứ đều rất bình thường, cho đến khi ai đó nhận ra một điều kỳ lạ.

Không ai khác, đó là Mio. Cô nhìn ngó quanh căn phòng, thắc mắc: ``Ủa? Kuon-senpai đâu rồi?''

Fubuki, người đang cầm lon nước ngọt, khựng lại một chút rồi liếc nhìn đồng hồ. Tám giờ ba mươi sáng. ``Hửm? Đã quá giờ vào làm rồi mà nhỉ\dots''

Mel nghiêng đầu, vẻ mặt ngơ ngác vì sự vắng mặt khó hiểu của cậu Tổng: ``Bình thường anh ấy nghỉ phép là sẽ thông báo trước mà, nhỉ?''

``Hay là nhà anh ấy có việc-nanodesu?'' Rushia lên tiếng, không khỏi băn khoăn.

Kanata khoanh tay, gật gù đồng tình: ``Anh ấy tự dưng nghỉ thế này đúng là lạ thật.''

Nene, người nãy giờ không theo kịp cuộc trò chuyện, giơ tay lên đầy bối rối: ``Ể, ể? Nene cần một ai đó giải thích-aru\dots''

Mọi người bắt đầu bàn tán xôn xao. Đối với họ, sự hiện diện của Kuon trong văn phòng gần như là chuyện hiển nhiên. Dù là Tổng quản lý, cậu vẫn thường xuyên hỗ trợ họ trong các buổi tập luyện, xử lý giấy tờ, hoặc đơn giản là có mặt để lắng nghe những câu chuyện của họ. Việc cậu không xuất hiện mà không báo trước khiến họ có cảm giác kỳ lạ, như thể một mảnh ghép quan trọng đã bị lấy mất.

Ngay lúc này, cửa văn phòng mở ra, A-chan bước vào. Vừa nhìn thấy nàng Đeo kính, các cô gái lập tức dồn sự chú ý về phía cô. Nhưng trước khi kịp nói gì, A-chan đã giơ tay ra hiệu: ``Trật tự nào! Mọi người nghe chị nói đã!''

Tuy vậy, cả căn phòng vẫn như một cái chợ vỡ. Họ sốt ruột tới mức tiếng người nọ chen ngang lời người kia, khiến cho cô nàng không biết phải xử trí thế nào.

Phải mất hai, ba lần hắng giọng, cả căn phòng mới dần yên lặng, dù vẫn có vài tiếng xì xào nho nhỏ. Đến khi cảm thấy mọi người đã sẵn sàng, nàng ``Tăng ca'' hít sâu một hơi và nói một câu ngắn gọn: ``Hôm nay công ty cho Hikawa-san nghỉ rồi.''

Một giây. Hai giây. Ba giây.

``\textbf{Ể!??}''

Cả căn phòng đồng loạt hét lên đầy kinh ngạc. Không chỉ có tiếng la hét, mà còn có tiếng ghế xê dịch, tiếng đồ đạc lạch cạch khi một số người giật nảy lên vì bất ngờ.

``K-Không thể như thế được!'' Noel ôm đầu, mặt tái mét lại.

Aqua thì hoang mang tột độ: ``T-Tại sao chứ!?''

Korone toát mồ hôi hột, sốt sắng hỏi lại A-chan: ``Anh ấy đã làm gì sai sao!?''

Mặt Lamy thì xanh xao như tàu lá chuối, khẽ rùng mình khi cô đã từng trải nghiệm một cảm giác mất mát như thế: ``Chắc phải có sự nhầm lẫn gì chứ!?''

Subaru thì tỏ vẻ quyết liệt: ``Anh ấy lại định trêu chúng ta nữa ư!?''

Sự lo lắng lan tỏa nhanh chóng. Họ không thể tin được Kuon lại bị công ty cho nghỉ mà không báo trước. Một vài người thậm chí đã nghĩ đến viễn cảnh tồi tệ nhất.

A-chan, dự đoán được phản ứng này từ trước, đành vỗ trán thở dài: ``Bình tĩnh nào, mấy đứa! Để chị giải thích đã!''

Nàng Sói đen, người vẫn giữ được sự điềm tĩnh, lên tiếng: ``Mọi người, cứ để cho A-chan nói gì đã nào. Chắc công ty phải có lý do gì thì mới để cho anh ấy nghỉ chứ, đúng không?''

``Đúng đó,'' nàng Đeo kính đáp. Cô chỉnh lại cặp kính, rồi tiếp tục thông báo: ``Thực ra, nghỉ ở đây ý chị không phải là bị đuổi việc.''

Những tiếng xì xào im bặt. Ba mươi mấy cặp mắt mở to nhìn nàng tóc tím như thể cô vừa nói một thứ gì đó khó tin.

``\dots Ể?'' / ``Hả?'' / ``Chị nói sao cơ?''

``Chị bảo `cho nghỉ' ở đây là nghỉ một ngày ấy,'' cô đáp, khẽ bật cười với phản ứng ngỡ ngàng của mọi người.

Sự căng thẳng trong không khí lập tức tan biến. Những tiếng thở phào nhẹ nhõm vang lên khắp phòng. Một vài người chống tay lên đầu gối, như thể vừa trút được một gánh nặng lớn.

``Chị lại dọa mọi người rồi, peko\dots'' Pekora thở phào nhẹ nhõm, đôi tai thỏ cụp xuống.

``Không vui tí nào đâu, A-chan à!'' Polka tức đến mức xì khói, nhìn cô nàng Đeo kính với ánh mắt không thể nào tin nổi.

A-chan ngừng cười, rồi nhún vai đáp lại: ``Mấy đứa có để chị nói hết đâu. Đến cả YAGOO còn không muốn đuổi cậu ta đi nữa là.''

Mọi người bật cười nhẹ. ``Phải rồi ha\dots'' / ``Làm gì có chuyện như thế!'' một số người trong nhóm thốt lên. Như vậy, Kuon chỉ nghỉ ở nhà, chứ không phải là bị nghỉ việc.

``Nhưng mà, nếu thế thì\dots'' Towa định đưa ra ý kiến gì đó.

Dường như tất cả đều đang nghĩ đến cùng một câu hỏi. Họ đồng loạt quay sang nhìn A-chan, đồng thanh: ``Tại sao hôm nay công ty lại cho anh ấy nghỉ?''

A-chan chỉ thở dài, khoanh tay lại rồi hỏi ngược lại: ``Mấy đứa biết hôm nay là ngày gì không?''

Một khoảnh khắc im lặng bao trùm căn phòng. Các cô gái đưa mắt nhìn nhau, cố gắng lục lọi ký ức xem ngày hôm nay có gì đặc biệt. Một số người bắt đầu đếm ngón tay, cố nhớ lịch trình công ty. Một số khác chỉ nghiêng đầu bối rối.

``Hôm nay là 17 tháng 7\dots'' Shion đáp lại.

``Ngoài việc hôm nay bọn em không có buổi stream nào ra,'' Marine tiếp lời, ``còn lại em không nghĩ ra được hôm nay là gì ạ.''

A-chan lắc đầu, như thể đoán định được mọi thứ diễn ra như vậy. Cô chậm rãi nói: ``Hôm nay là sinh nhật của Hikawa-san đó.''

``\dots Ể?''

Phản ứng đầu tiên đến từ những người cuối cùng cũng nhớ ra. Các thành viên cộm cán nhất của \hll\ - từ Thế hệ thứ Không và số ít thuộc Thế hệ thứ Nhất, thứ Hai - có lẽ là những người nhớ ra được.

``Chẳng trách mà em cảm thấy hôm nay là ngày gì trọng đại lắm\dots'' Sora nhỏ nhẹ đáp.

Suisei đập trán, thở dài: ``Sao chúng ta có thể quên được nhỉ?''

Shion chớp mắt vài lần, rồi khẽ gật gù: ``Bảo sao em nghe thấy quen thế\dots''

Nhưng không phải ai cũng có phản ứng như vậy - thực tế là phần lớn các thành viên còn lại. Một số người chỉ biết đứng đó, hoàn toàn ngơ ngác.

``Bây giờ em mới biết đấy,'' Botan nghiêng đầu.

Ayame khoanh tay, khẽ cau mày: ``Anh ấy có nói cho chúng ta biết đâu nhỉ?''

``Nghe khá nà kỳ đó na\~{}'' Luna ngân nga, vỗ tay.

Fubuki, người vẫn đang xoay lon nước ngọt trong tay, nghiêng đầu khó hiểu: ``Nhưng nếu thế thì anh ấy được cho nghỉ để làm gì?''

A-chan mỉm cười, ánh mắt ánh lên chút tinh nghịch: ``Mừng vì em đã hỏi, Shirakami-san. Chúng ta sẽ tạo bất ngờ cho cậu ấy.''

Một làn sóng xôn xao lập tức lan khắp căn phòng. Các thành viên nhìn nhau, thì thầm bàn tán.

``Bất ngờ kiểu gì?''

``Ai sẽ lo phần gì đây?''

``Liệu có nên rủ cả gia đình anh ấy không nhỉ?''

Một kế hoạch bí mật bắt đầu thành hình.

\ct

Ngay sau khi A-chan đưa ra ý tưởng, cả căn phòng rơi vào một trận bàn tán sôi nổi.

Subaru vỗ hai tay vào nhau, mắt sáng rực: ``Nếu đã làm thì phải làm cho đàng hoàng! Sinh nhật mà không có tiệc thì còn gì là sinh nhật nữa, đúng không!?''

``Đồng ý, ne!'' - Miko gật đầu. ``Nhưng mà làm tiệc ở đâu bây giờ?''

``Nhà anh ấy chứ đâu!'' Nene hào hứng giơ tay. ``Chúng ta có thể xông vào nhà anh ấy rồi hét `Ngạc nhiên chưa\~{}!' như trong phim ấy-aru!'

Lamy kéo tay cô bạn lại, cười trừ: ``Khoan đã, Nene-chi. Nếu không báo trước mà cứ xông như thế thì chắc anh ấy sẽ choáng mất\dots''

``Vậy chúng ta nên hỏi gia đình Kuon-senpai trước nhỉ?'' Flare trầm ngâm.

Tất cả mọi người quay về phía A-chan với những cặp mắt mong chờ.

Thấy vậy, nàng Đeo kính gật đầu, xác nhận: ``Chị đã nghĩ tới tình huống đó rồi, nên chị đã liên hệ trước với gia đình Hikawa. Họ rất sẵn lòng hợp tác với chúng ta.''

``Ô\~{}, vậy thì dễ rồi!'' Korone vui vẻ nắm tay cô bạn thân của mình lại. ``Mọi người sẽ đến nhà Kuon-senpai, còn mình sẽ lo bánh kem!''

``Hử? Mà mấy đứa đã ai từng thấy Kuon-senpai ăn bánh chưa?'' Pekora chống tay lên cằm, thắc mắc.

Cả phòng bỗng chốc im bặt. Một số người lẩm bẩm suy nghĩ, cố nhớ lại xem đã bao giờ thấy cậu Tổng của họ ăn bánh hay chưa.

``Ể? Anh ấy có ăn đồ ngọt không nhỉ?'' Watame nghiêng đầu, thoáng bối rối. Dường như không ai nghĩ tới trường hợp này.

Towa khoanh tay, nhíu mày: ``Anh ấy từng nói thích uống trà, nhưng chưa bao giờ nhắc đến bánh ngọt\dots''

``Chà\dots\ khó đây,'' Botan bật cười. ``Nếu làm bánh kem mà chính chủ không ăn thì cũng hơi buồn nhỉ?''

``Vậy thử hỏi gia đình anh ấy xem sao, ne?'' Miko đề xuất.

Okayu gật đầu đáp lại: ``Ý hay á\~{} Bọn mình sẽ hỏi bố mẹ hoặc em gái anh ấy\~{}''

``Ừm, mà nói mới nhớ\dots'' Shion khoanh tay lại. ``Quà thì sao? Chúng ta cần mua quà nữa mà.''

Một lần nữa, cả căn phòng rơi vào trầm tư. Mua quà cho Kuon? Nhưng nên mua cái gì đây?

Flare chống cằm, suy nghĩ: ``Anh ấy không có sở thích rõ ràng lắm, nhỉ?''

``Chưa từng thấy anh ấy nói muốn mua gì cả, dazo\dots'' Ayame đồng tình. ``Ca này có vẻ hơi khó đấy.''

A-chan bật cười trước sự bối rối của các cô gái: ``Chắc chúng ta cần một buổi điều tra sở thích của cậu ta thôi nhỉ?''

Tất cả đồng loạt gật đầu. Một kế hoạch tỉ mỉ bắt đầu hình thành trong đầu họ.

Nàng Đeo kính liền lấy điện thoại ra và gọi cho Ryujin. Để đảm bảo rằng tất cả mọi người nghe được, cô đã chủ động bật loa ngoài lên.

Cô nàng là một trong số những người hiếm hoi ngoài những khách hàng của mình mà vẫn giữ số liên lạc thường xuyên với cậu Tổng.

Điện thoại vừa đổ chuông hai tiếng đã có người nhấc máy. Một giọng trầm đặc vang lên ở đầu dây bên kia, giọng hơi luyến một chút: ``Alo?''

``Chú Ryujin-san, chào chú! Cháu gọi để hỏi một chút về con nhà chú ạ.''

``Ồ? Thằng nhóc gây rắc rối gì à?''

A-chan bật cười: ``Không đâu ạ. Hôm nay là sinh nhật của Hikawa-san, nên tụi cháu đang tìm hiểu một chút về sở thích của cậu ấy để mua quà cho phù hợp ạ.''

``À, ra vậy\dots\ Chờ chú một lát.''

Ryujin hắng giọng, rồi lặng lẽ bước xuống tầng một. Trước đó, ông đã nói với Kuon rằng có khách hàng đến nên sẽ bận một lúc. Đảm bảo rằng con trai không nghe thấy cuộc gọi, ông mới tiếp tục: ``Thằng bé à\dots\ Bình thường lúc rảnh nó thường nghe nhạc và chơi game.''

``Game thì tụi cháu cũng đoán được, nhưng nghe nhạc à?'' nàng Đeo kính tò mò.

``Ừ. Nó thường nghe nhạc của các thành viên \hll, hình như là để kiểm tra chất lượng.''

Cả căn phòng bất giác rộn lên.

``Ể\~{}? Anh ấy có nghe nhạc của bọn em thật sao!?'' nàng Cáo trắng phấn khích.

Pekora tròn mắt: ``Sao chưa bao giờ thấy anh ấy nhắc đến chứ, peko!?''

``Anh ấy giấu kỹ quá nhỉ\dots'' Towa lẩm bẩm, không tin nổi vào tai mình.

Nhưng ngay sau đó, người bố có bổ sung: ``Ngoài ra nó cũng hay nghe nhạc từ game nữa.''

Bầu không khí vui vẻ trong phòng bỗng nhiên trở nên ngơ ngác.

``\dots Hả?'' nàng Vịt chớp mắt. ``Ý chú là sao ạ?''

``Nhạc trong game ấy. Những bài nhạc nền, nhạc chủ đề, đại loại thế. Nó còn có cả playlist (danh sách nhạc) riêng luôn.''

Toàn bộ phòng họp chìm vào im lặng vài giây. Họ không nghĩ cậu Tổng của họ có sở thích kỳ quặc như vậy.

``\dots Anh ấy nghiêm túc chứ?'' Aqua lên tiếng, giọng như không tin vào tai mình.

``Anh ấy nghe nhạc game để giải trí à?'' Ayame bối rối.

`` Không phải là mấy bài hát có ca sĩ hẳn hoi đâu nhỉ?'' Mel nghiêng đầu thắc mắc.

Ryujin bật cười: ``Không, nó nghe cả mấy bài không lời luôn. Có lần chú còn thấy nó nghe đi nghe lại một bản nhạc nền của một game cũ rích nào đấy, mà nó nói là `hay lắm.' Chú cũng không hiểu tại sao luôn.''

``\dots Cái gu này thì bọn cháu cũng không ngờ tới thật\dots'' Roboco cười trừ.

Dù hơi bất ngờ, các thành viên \hll\ cũng dần tiếp nhận được thông tin này. Có vẻ như họ sẽ phải suy nghĩ kỹ hơn về quà tặng dành cho Kuon.

``Cảm ơn chú nhiều ạ!''

``Không có gì. Nhớ giữ bí mật nhé, thằng bé mà biết thì hỏng chuyện đấy.''

``Bọn cháu hiểu rồi!''

A-chan cúp máy, rồi nhìn quanh phòng. ``Rồi, bước đầu coi như xong. Tiếp theo sẽ là\dots\ đồ ăn.''

``Hmmm\dots\ vậy rốt cuộc anh ấy thích ăn gì nhỉ?'' Nene nghiêng đầu suy nghĩ.

``Có ai từng thấy Kuon ăn gì nhiều chưa?'' Noel lên tiếng.

``Toàn thấy anh ấy ăn suất văn phòng bình thường thôi. Không thì cũng cơm nắm hoặc là sushi, hoặc là bánh mì à,'' Mio nhún vai đáp, hầu như không có ý tưởng gì.

``Không chừng senpai thuộc kiểu dễ tính, cho gì ăn nấy ấy\dots'' Marine phỏng đoán.

A-chan sau khi tiếp thu ý kiến từ ba người, gật đầu, nhưng vẫn tỏ vẻ không chắc chắn. ``Để chắc chắn thì ta gọi cho mẹ cậu ấy đi.''

Lần này, Sora là người nhận trách nhiệm. Cô lấy điện thoại ra, bấm số và gọi cho Kaien. Khác với người chồng, người phụ nữ này có thể nhận bất cứ cuộc gọi nào, thậm chí còn bắt cả số của\dots\ bên lừa đảo, mặc dù bà biết rất rõ.

Chỉ sau hai hồi chuông, giọng nói dịu dàng của người phụ nữ vang lên: ``Sora-chan đó sao?''

``Cháu chào cô Kaien-san ạ!''

``Chào cháu! Có chuyện gì vậy?''

``Dạ, bọn cháu đang chuẩn bị cho sinh nhật Kuon-senpai, nên muốn hỏi về món ăn mà anh ấy thích nhất ạ.''

Kaien đang ở trong bếp, bà tranh thủ đi ra chỗ khác để đảm bảo Kuon không nghe thấy. ``Ồ, ra là vậy. Hmm\dots\ thật ra thì thằng bé khá dễ nuôi, cho gì cũng ăn.''

``\dots Ể?'' Sora chớp mắt. Lại một điều nữa mà cô không ngờ tới.

``Anh ấy ăn được hết á?'' Fubki trố mắt ngạc nhiên.

``Không kén ăn tí nào sao?'' Lmay nhướng mày, hỏi lại người mẹ của cậu Tổng.

Bà ấy im lặng một lúc lâu, như thể ngẫm lại sở thích của cậu quý tử nhà mình. ``\dots Cũng không hẳn. Nó không thích ăn rau lắm đâu.''

``\textbf{Ể!!?}'' một người trong nhóm hét lên đầy vui sướng.

Tất cả quay sang nhìn Suisei, người lúc này đang cười rạng rỡ như thể vừa tìm thấy tri kỷ. ``Mình biết mà! Một người mạnh mẽ như Kuon-senpai thì sao có thể thích rau củ chứ!''

Một số người trong nhóm nở những nụ cười trừ tới cô nàng Ca sĩ vui vẻ kia. A-chan thì chỉ có thể lắc đầu: ``Đúng kiểu Hoshimachi-san luôn\dots''

Bỏ qua sự phấn khích quá đà của nàng Ca sĩ, Kaien tiếp tục: ``Nhưng nếu nói món nó thích nhất\dots\ thì chắc là thịt bò.''

``\textbf{Thịt bò!?}'' một giọng nói phấn khích khác vang lên.

Lần này là Noel. Cô nắm chặt tay, đôi mắt sáng rực như trúng số độc đắc: ``Thế thì hợp với chị quá còn gì!''

Thậm chí, cô nàng còn tới chung vui với Suisei - cũng là người thích ăn các loại thịt, càng làm cho các cô gái còn lại toát mồ hôi hột.

``Vậy là Kuon-senpai thích thịt bò à\dots'' Mio gật gù.

``À mà, nó cũng ăn được đồ ngọt nữa. Có lần sinh nhật cậu nó hôm bữa, nó đánh chén hết một phần tư cái bánh luôn à!''

``Vậy là làm bánh kem được rồi!'' Korone hào hứng nói, đuôi vẫy vẫy mạnh.

Sora cười nhẹ rồi nói vào điện thoại: ``Dạ, vậy là tụi cháu nắm được rồi ạ. Cảm ơn cô Kaien nhiều lắm!''

``Không có gì đâu. Nhớ giữ bí mật đấy nhé!''

``Dạ, bọn cháu sẽ cẩn thận ạ!''

Sau khi cuộc gọi kết thúc, A-chan gõ tay lên bàn. ``Rồi, vậy là chúng ta có thông tin về sở thích và đồ ăn của Kuon rồi. Giờ còn một chuyện quan trọng nữa\dots''

``Quà tặng đúng không?'' Roboco mỉm cười.''

``Đúng thế. Để chắc ăn, ta nên hỏi người hiểu Kuon nhất\dots''

Mọi ánh mắt đổ dồn về phía cô. A-chan nở một nụ cười. ``Ta sẽ hỏi em gái của cậu ấy.''

Nàng Sói đen mỉm cười, lấy điện thoại ra gọi cho cô em gái của Kuon.

Cô bấm số và đưa lên tai. Sau vài giây chờ đợi, giọng của một cô gái trẻ vang lên ở đầu dây bên kia, hòa với chút tạp âm khác.

``A-lô? Mio-onee-chan ạ?''

``Chị đây. Ryuuhou-chan đang ở trường hả?''

``Dạ, vẫn đang trong giờ nghỉ. Có chuyện gì không chị?''

``Chị muốn hỏi một chút\dots về sinh nhật của Kuon-senpai.''

``À, ra thế\dots\ Onii-san vẫn không nói gì với mọi người đúng không?''

``Ừ\dots'' không chỉ nàng Sói, mà tất cả mọi người cùng đồng thanh lên tiếng, tỏ vẻ hơi bất mãn.

Họ nghe thấy tiếng thở dài từ đầu dây bên kia, như thể bản thân cô em gái của cậu cũng thấy bất bình thay cho họ: ``Haiz\dots\ Đúng là kiểu của anh ấy.''

Nàng Sói đen hắng giọng, rồi đi thẳng vào vấn đề: ``Thế nên bọn chị muốn chuẩn bị quà cho bất ngờ. Em có biết anh ấy thích quà gì không?''

Ryuuhou im lặng một chút, rồi mới ngập ngừng trả lời: ``Nói thật thì\dots\ anh ấy không phải kiểu người thích nhận quà.''

``\textbf{Ể!?}''

Những người xung quanh Mio đều chớp mắt ngạc nhiên.

``K-Không thích nhận quà?'' Polka lặp lại.

``Nhưng nếu có ai tặng thì chắc là anh ấy vẫn nhận chứ?'' Subaru hỏi tường tận.

Cô em gái thở dài một tiếng, rồi tiếp tục: ``Vẫn nhận ạ. Nhưng không phải vì ham thích mà vì phép lịch sự thôi. Onii-san không phải là người ham vật chất mà.''

Mio nhíu mày: ``Vậy thì khó nghĩ ghê\dots\ Không lẽ bọn chị không cần tặng gì hết?''

Cô nhóc ở bên kia cười nhẹ. ``Không hẳn đâu. Thật ra, có một lần anh ấy từng nói với em rằng\dots\ nếu có ai muốn tặng quà, thì anh ấy cũng thích có một cái máy chơi game cầm tay, như S*itch chẳng hạn.''

``Hả!?'' một số người thốt lên. Đây là một trong những thiết bị chơi điện tử mà các thành viên \hll\ hay dùng khi có buổi trực tiếp với nhau.

Nhưng việc mà cậu Tổng - một người làm việc đến bù đầu, kiếm rất nhiều tiền - mà lại không có nó, lại là một điều kỳ quặc trong con mắt của họ.

``Nhưng mà\dots\ anh ấy có tiền mà?'' Mio ngạc nhiên hỏi. ``Tại sao anh ấy lại không tự tậu cho mình một chiếc chứ?''

``Thì đúng là thế,'' Ryuuhou nói, ``Chẳng qua Onii-san anh ấy không mua vì thấy nó không thiết thực lắm thôi.''

``Kiểu\dots\ không đến mức cần thiết ấy hả?''

``Dạ, kiểu vậy đó. Anh ấy có nói là chơi game trên điện thoại hay trên máy tính là cũng ổn rồi,'' cô đáp lại. ``Nhưng nếu có ai tặng thì chắc chắn anh ấy vẫn sẽ dùng thôi.''

Mọi người xung quanh Mio gật gù với đề xuất này.

``Ừm\dots\ nếu là S*itch thì có vẻ hợp lý,'' Fubuki suy nghĩ. ``Chúng ta nhiều khi cũng muốn rủ anh ấy chơi cùng lắm chứ.''

``Nếu là game, thì chắc chắn Kuon-senpai sẽ thích rồi,'' Towa nhận xét.

``Chúng ta có thể góp tiền mua một bộ đầy đủ, cả máy lẫn game luôn,'' Aki đề xuất.

``Một ý tưởng không tồi,'' A-chan khoanh tay, mỉm cười gật đầu.

Nàng Sói đen quay lại điện thoại. ``Rồi, cảm ơn Ryuuhou-chan nhé! Em giúp bọn chị nhiều lắm!''

Cô em gái cười khì: ``Không có gì đâu ạ. Nhớ giữ bí mật nha!''

``Tất nhiên rồi!'' cả ba mươi mấy người đồng thanh.

Sau khi kết thúc cuộc gọi, Mio nhìn quanh tất cả mọi người, như thể họ đã biết được sở thích của cậu Tổng: ``Vậy là xong rồi nhỉ?''

A-chan khoanh tay: ``Vậy là bây giờ chúng ta đã có đủ thông tin rồi. Sở thích, đồ ăn và quà sinh nhật. Kế hoạch tiếp theo sẽ là\dots''

``\dots mua sắm!'' tất cả bọn họ reo lên.

Cả căn phòng lập tức trở nên náo nhiệt.

``Đi lựa quà nào!''

``Chúng ta nên phân nhóm ra để dễ tìm!''

``Còn bánh sinh nhật thì sao? Ai sẽ phụ trách?''

``Có ai muốn đặt đồ handmade không!?''

Nàng Đeo kính thở dài, nhìn những cô gái đang tràn đầy năng lượng. Cô mỉm cười.

``Được rồi. Hãy cùng nhau làm cho sinh nhật lần này của Kuon-senpai trở thành một kỷ niệm đáng nhớ nào!'' nàng Thần tượng làm nóng tinh thần của mọi người.

Tất cả đồng loạt giơ tay hô lớn: ``\textbf{OHHHHH!!!}''

\begin{center}
    \uline{Vào cùng thời điểm\\Nhà Hikawa.}
\end{center}

\pov{Hikawa Kuon}

Hôm nay là ngày Mười Bảy tháng Bảy.

Nói đúng hơn, hôm nay là sinh nhật của tôi.

Thêm một tuổi mới. Cũng như mọi năm, tôi chẳng cảm thấy có gì đặc biệt cả. Sinh nhật vốn chỉ là một dấu mốc cho thấy thời gian đang trôi đi, và với tôi, nó cũng chỉ là một ngày bình thường như bao ngày khác. Không có gì đáng để háo hức hay mong chờ giống như hồi còn bé cả.

Năm nào cũng vậy, gia đình tôi sẽ ăn một bữa tối đặc biệt hơn ngày thường một chút, có thể mua một cái bánh kem nho nhỏ để thêm không khí, rồi kết thúc trong yên bình. Không tiệc tùng rầm rộ, không tổ chức linh đình. Tôi cũng chẳng mấy khi nói với ai về ngày sinh của tôi, đơn giản vì không muốn làm phiền người khác. Nếu ai đó nhớ thì tốt, còn không thì cũng chẳng sao.

Tuy nhiên, có một điều hơi khác so với mọi năm: cả bố mẹ tôi đều đang có mặt ở nhà, và dường như họ đang bàn bạc điều gì đó. Ryuuhou thì vẫn đang đi học, nên trong nhà lúc này chỉ có ba người chúng tôi.

Ngoại trừ việc hôm nay là sinh nhật tôi mà bố tôi vẫn còn chút việc từ khách hàng. Biết sao được, hôm nay có phải là ngày nghỉ đâu. Tôi cũng không có ý kiến gì với việc đó cả.

Tôi ở ngồi bàn, chống cằm, nhìn bố mẹ và hỏi những câu hỏi mà năm nào cũng vậy: ``\vie{Năm nay chắc lại ăn sinh nhật như mọi năm thôi nhỉ?}''

Bố tôi nghe vậy, nhướng mày: ``\vie{Cả nhà mình ăn với nhau à?}''

``\vie{Vâng. Con thấy chỉ cần cả nhà ăn là cũng vui rồi. Coco-chan với Aloe-chan lại có việc riêng nữa, nên chắc họ không ở nhà đâu.}''

Thì, đúng là thế. Nàng Rồng thì đang đi bốc vác ở kho mẹ tôi, còn nàng Succubus thì đang nhận một show diễn ở tận miền Trung cơ mà.

Mẹ tôi, người vừa từ trên bếp xuống, nghe tôi nói vậy, có vẻ hơi ngạc nhiên: ``\vie{Không định rủ thêm bạn đến à?}''

``\vie{Chúng nó cũng có việc chứ. Có phải ai cũng rảnh như con đâu.}''

``\vie{Thế còn các cô gái ở công ty thì sao?}'' bố tôi hỏi.

``\vie{Họ cũng có việc của họ mà. Mình con nghỉ ở nhà. Thế mới lạ!}''

Tới đây, bất chợt mẹ tôi trố mắt, chỉ ra một điểm hơi bất thường: ``\vie{Ơ mà, hôm nay là thứ bảy mà. Mẹ tưởng hôm nay con đi làm chứ?}''

Tôi thở dài một hơi, gõ nhẹ ngón tay xuống mặt bàn: ``\vie{Thì định thế, nhưng tối qua A-san bảo hôm nay con nghỉ.}''

Bố tôi nhìn tôi với ánh mắt tò mò: ``\vie{Tại sao lại như thế?}''

Tôi thở dài một hơi. ``\vie{Chuyện là thế này\dots}''

\begin{center}
    \uline{Hồi tưởng\\Tối hôm qua.}
\end{center}

Lúc đó đã là chín giờ tối.

Tôi vẫn còn đang cặm cụi kiểm tra lại bản báo cáo cần nộp vào tuần sau. Cả ngày đã bận rộn với công việc, nhưng dù sao cũng phải đảm bảo không có sai sót. Khi đang lướt qua những dòng cuối cùng, bỗng nhiên điện thoại của tôi đổ chuông.

Người gọi đến, không ai khác ngoài A-san ra.

Tôi nhấn nghe máy: ``Vâng, nhà Hikawa xin nghe.''

Giọng của chị ấy vang lên: ``Hikawa-san đúng không? Tôi có việc này muốn nói với cậu đây.''

``Vâng, chị cứ nói.''

Việc chị ấy gọi cho tôi vào giờ này cũng không phải điều lạ lẫm gì. Thậm chí, có lúc A-san còn gọi cho tôi lúc mười hai giờ đêm nữa, nên tôi đã quá quen với những cuộc gọi như thế này.

Nhưng\dots\ điều mà chị ấy sắp chuyển cho tôi làm cho tôi cảm thấy khó hiểu.

``Giám đốc chuyển lời cho cậu bảo ngày mai cậu không phải đi làm nữa nha.''

Tôi hơi khựng lại, mắt chớp chớp vài cái: ``Ơ? Là sao? Em bị đuổi-''

``Để tôi nói hết đã nào,'' A-san cắt ngang ngay lập tức. ``Nghỉ ở đây là mai nghỉ một ngày ấy. Chúng ta chưa có ý định đuổi ai đi kể từ khi Kiryu-san tốt nghiệp đâu.''

Thấy thế, tôi thở phào nhẹ nhõm. Không biết có phải do nghiệp quật từ khi tôi dọa Suba-chan vụ tôi nộp đơn nghỉ không, nhưng pha đó làm cho tôi hú hồn.

Tuy thế, tôi thấy vẫn còn ngơ ngác trước yêu cầu đầy nghi vấn này: ``Nhưng mà tại sao em được cho nghỉ ạ?''

A-chan im lặng một chút, rồi hạ giọng: ``Cậu biết rõ ngày mai là ngày gì mà.''

Tôi nhíu mày, suy nghĩ một lúc. Rồi\dots

``À\dots\ Phải. Ngày mai là sinh nhật em.''

Tôi khẽ gật đầu với suy nghĩ này. Nhưng điều đó lại làm cho tôi thêm thắc mắc: ``Cái đó thì\dots\ liên quan gì tới việc YAGOO-san lại cho em nghỉ ạ?''

Đầu dây bên kia vang lên một tiếng thở dài: ``Tôi cũng không biết đâu. YAGOO-sama nói thế thì tôi cũng chuyển lại lời cho cậu thôi. Chẳng biết ông ta đang nghĩ gì nữa.''

Tôi cũng chỉ biết đường thở dài theo. ``Vâng, biết thế\dots\ Vậy em sẽ không hỏi chị nữa.''

``Hình như ông ấy nói là sẽ tạo bất ngờ gì đó\dots\ Tôi không biết được, nhưng cậu cứ chuẩn bị đi cho chắc.''

Tạo bất ngờ?

Nghe xong câu đó, tôi cảm thấy có chút kỳ lạ. Nhưng tôi cũng không nghĩ quá nhiều về nó. Tôi chỉ lịch sự đáp gọn: ``Dạ, em biết rồi.''

\begin{center}
    \uline{Kết thúc hồi tưởng.}
\end{center}

``\vie{\dots Và chuyện là thế đấy ạ.}''

Sau khi nghe xong câu chuyện của tôi, bố tôi khoanh tay, gật gù: ``\vie{À\dots{} Tức là giám đốc bảo con nghỉ nhưng không nói lý do chứ gì?}''

Tôi nhún vai đáp lại: ``\vie{Vâng\dots\ Con cũng chẳng biết vì sao nữa. Bảo là bất ngờ gì ấy\dots}''

Tôi thở dài một hơi. Dù gì thì tôi cũng không có kế hoạch gì cho ngày hôm nay, nên được nghỉ cũng không tệ lắm. Nhưng tôi vẫn không hiểu YAGOO-san đang định làm gì.

``\vie{Thôi kệ đi. Dù sao thì, hôm nay chắc vẫn như mọi năm thôi, con không nghĩ có gì đặc biệt xảy ra đâu.}''

Mẹ tôi gật đầu: ``Vậy nhà mình sẽ ăn như mọi năm nhỉ?''

Tôi nhún vai đáp lại đầy tỉnh bơ: ``Lẩu hoặc nướng, dạng kiểu thế ạ.''

Trộm vía tôi không phải người quá kén ăn, nên món gì cũng được, miễn là cả nhà có thể quây quần bên nhau.

Cả nhà gật gù đồng tình, rồi bắt đầu chuẩn bị đi chợ, lo nguyên liệu cho buổi sinh nhật nho nhỏ tối nay. Mọi chuyện vẫn bình thường như mọi năm. Chẳng có gì bất ngờ cả.

Thực ra, mong ước của tôi cũng đơn giản thôi: có một bữa tiệc sinh nhật yên bình. Không cần phải cầu kỳ hay hoành tráng, chỉ cần một bữa ăn ấm cúng cùng gia đình là đủ.

Dù vậy\dots

Trong thoáng chốc, tôi cũng đã nghĩ đến một viễn cảnh khác. Không biết có phải là tươi sáng hơn hay u tối hơn.

Giả sử các thành viên \hll\ bất ngờ kéo đến nhà tôi, chúc mừng sinh nhật, thì sao nhỉ?

Nó sẽ vui đấy. Nhưng mà\dots\ thật sự thì cũng khá viển vông. Ai cũng có công việc riêng, ai cũng có lịch trình dày đặc. Họ có thể nhớ ngày sinh của tôi, nhưng điều đó không có nghĩa là họ sẽ đến tận nơi chỉ để tổ chức sinh nhật cho tôi.

Tôi cười nhạt, phẩy tay với suy nghĩ điên rồ đó.

Không sao. Không đòi hỏi quá nhiều là được.

Nhưng\dots\ liệu có thực sự như vậy không?

\end{document}